\subsection{Quadratic Forms and Polar Forms}

\begin{definition} \label{def:quadratic-form-and-polar-form} \hlt{Quadratic Forms and Polar Forms}\\
A \hlt{quadratic form} on $\mathcal{V}$ is a function $Q:\mathcal{V}\to\mathbb{F}$ such that:
\begin{enumerate}[label=(\roman*)]
\setlength{\itemsep}{0pt}
\item $Q(a\mathbf{v})=a^2Q(\mathbf{v})$ for all $a\in\mathbb{F}$ and $\mathbf{v}\in\mathcal{V}$;
\item the \hlt{polar form} $B_Q(\mathbf{u},\mathbf{v})\defeq Q(\mathbf{u}+\mathbf{v})-Q(\mathbf{u})-Q(\mathbf{v})$ is bilinear.
\end{enumerate}
\end{definition}

\begin{proposition} \label{prop:quadratic-forms-symmetric-bilinear-forms} \hlt{Quadratic Forms and Symmetric Bilinear Forms}\\
For every quadratic form $Q$, its polar form $B_Q$ is symmetric.\\
If $\operatorname{char}(\mathbb{F})\neq2$, maps $\{B:\operatorname{Sym}(B)\} \underset{Q\mapsto B_Q}{\overset{B\mapsto Q_B}{\rightleftarrows}} \{Q:Q\text{ quadratic}\}$ and $Q_B(\mathbf{x})\defeq\frac{1}{2}B(\mathbf{x},\mathbf{x})$ are inverse correspondences between symmetric bilinear forms and quadratic forms.
\begin{proof} First, $B_Q(\mathbf{u},\mathbf{v}) =Q(\mathbf{u}+\mathbf{v})-Q(\mathbf{u})-Q(\mathbf{v}) =B_Q(\mathbf{v},\mathbf{u})$, so $B_Q$ is symmetric.
\begin{enumerate}[label=\roman*.]
\setlength{\itemsep}{0pt}
\item[$(\Rightarrow)$] If $\operatorname{Sym}(B)$, then $Q_B(a\mathbf{x})=a^2Q_B(\mathbf{x})$ and $B_{Q_B}(\mathbf{u},\mathbf{v}) =\frac{1}{2}\bigl(B(\mathbf{u}+\mathbf{v},\mathbf{u}+\mathbf{v})-B(\mathbf{u},\mathbf{u})-B(\mathbf{v},\mathbf{v})\bigr) =B(\mathbf{u},\mathbf{v})$.
\item[$(\Leftarrow)$] If $Q$ is quadratic, then $Q_{B_Q}(\mathbf{x}) =\frac{1}{2}B_Q(\mathbf{x},\mathbf{x}) =\frac{1}{2}\bigl(Q(2\mathbf{x})-2Q(\mathbf{x})\bigr) =Q(\mathbf{x})$.
\end{enumerate}
\hlt{Structural Commentary}: Away from characteristic $2$, quadratic and symmetric bilinear data coincide.\\
\hlt{Key Insight}: Dividing by $2$ recovers diagonal data from polar data.
\end{proof}
\end{proposition}

\begin{proposition} \label{prop:polar-forms-alternating-in-characteristic-two} \hlt{Polar Forms in Characteristic Two}\\
If $\operatorname{char}(\mathbb{F})=2$, then $\operatorname{Alt}(B_Q)$ for every quadratic form $Q$ on $\mathcal{V}$.
\begin{proof} For $\mathbf{x}\in\mathcal{V}$, $B_Q(\mathbf{x},\mathbf{x}) =Q(2\mathbf{x})-Q(\mathbf{x})-Q(\mathbf{x}) =Q(\mathbf{0})=0$. Thus $\operatorname{Alt}(B_Q)$.\\
\hlt{Structural Commentary}: In characteristic $2$, every polar form is alternating.\\
\hlt{Key Insight}: The diagonal of the polar form vanishes automatically because both $2\mathbf{x}$ and $2Q(\mathbf{x})$ vanish.
\end{proof}
\end{proposition}

\begin{example} \label{ex:characteristic-two-polar-form-loses-quadratic-data} \hlt{Characteristic Two Loses Diagonal Data}\\
Assume $\operatorname{char}(\mathbb{F})=2$. On $\mathbb{F}^2$, let $Q_0(x,y)=xy$ and $Q_1(x,y)=xy+x^2$.\\
Then $B_{Q_0}((x,y),(u,v))=xv+uy=B_{Q_1}((x,y),(u,v))$, and $Q_0(1,0)=0\neq1=Q_1(1,0)$.\\
Indeed, the added square has zero polar contribution, where $(x+u)^2-x^2-u^2=2xu=0$.\\
More generally, for any linear $L:\mathcal{V}\to\mathbb{F}$, $Q$ and $Q+L^2$ have the same polar form.\\
\hlt{Structural Commentary}: In characteristic $2$, squares are invisible to polarization.\\
\hlt{Key Insight}: Characteristic $2$ quadratic forms need data beyond their polar forms.
\end{example}

\begin{example} \label{ex:characteristic-two-alternating-versus-nonalternating-symmetric} \hlt{Alternating versus Nonalternating Symmetric Forms}\\
Assume $\operatorname{char}(\mathbb{F})=2$.\\
The symmetric matrix $H=\begin{bmatrix}0&1\\1&0\end{bmatrix}$ defines a nonsingular alternating form on $\mathbb{F}^2$, since $\begin{bmatrix}x&y\end{bmatrix}H\begin{bmatrix}x\\y\end{bmatrix}=2xy=0$.\\
Thus $H$ is not congruent to any nonzero diagonal form: alternating means all self-pairings vanish, while a nonzero diagonal form has some basis vector with nonzero self-pairing. By contrast, $[1]$ is symmetric and nonalternating.\\
\hlt{Structural Commentary}: In characteristic $2$, symmetric forms split into alternating and nonalternating parts.\\
\hlt{Key Insight}: Nonsingular alternating blocks are invisible on the diagonal.
\end{example}

\begin{proposition} \label{prop:orthogonal-isometry-preserves-quadratic-form} \hlt{Isometries from Quadratic Values}\\
Let $\mathcal{V}$ be an orthogonal geometry with $\operatorname{char}(\mathbb{F})\neq2$, let $q(\mathbf{x})\defeq\innerprod{\mathbf{x},\mathbf{x}}$, and let $T:\mathcal{V}\to\mathcal{V}$ be bijective. Then $T$ is an isometry if and only if $q(T\mathbf{x})=q(\mathbf{x})$ for all $\mathbf{x}\in\mathcal{V}$.
\begin{proof}
\begin{enumerate}[label=\roman*.]
\setlength{\itemsep}{0pt}
\item[$(\Rightarrow)$] If $T$ is an isometry, then $q(T\mathbf{x})=\innerprod{T\mathbf{x},T\mathbf{x}}=\innerprod{\mathbf{x},\mathbf{x}}=q(\mathbf{x})$.
\item[$(\Leftarrow)$] If $q\circ T=q$, then $\innerprod{T\mathbf{x},T\mathbf{y}} =\frac{q(T\mathbf{x}+T\mathbf{y})-q(T\mathbf{x})-q(T\mathbf{y})}{2} =\frac{q(\mathbf{x}+\mathbf{y})-q(\mathbf{x})-q(\mathbf{y})}{2} =\innerprod{\mathbf{x},\mathbf{y}}$.\\
Thus $T$ preserves the form, so it is an isometry.
\end{enumerate}
\hlt{Structural Commentary}: In odd characteristic, $q$ determines the orthogonal form.\\
\hlt{Key Insight}: Polarization turns preservation of lengths into preservation of inner products.
\end{proof}
\end{proposition}

\begin{definition} \label{def:metric-orthogonality-isotropy-radical} \hlt{Orthogonality, Isotropy, and Radical}\\
Let $\mathcal{V}$ be a metric vector space.
\begin{enumerate}[label=(\roman*)]
\setlength{\itemsep}{0pt}
\item $\mathbf{x}\perp\mathbf{y}$ means $\innerprod{\mathbf{x},\mathbf{y}}=0$; then $\mathbf{x},\mathbf{y}$ are \hlt{orthogonal}. For $S\subseteq\mathcal{V}$, $S^{\perp}\defeq\{\mathbf{v}\in\mathcal{V}:\innerprod{\mathbf{v},\mathbf{s}}=0\ \forall\,\mathbf{s}\in S\}$
\item $\mathcal{S}\leq\mathcal{V}$ is \hlt{totally isotropic} or \hlt{totally degenerate} if $\mathcal{S}\subseteq\mathcal{S}^{\perp}$, and \hlt{nonsingular as a subspace} if $\mathcal{S}\cap\mathcal{S}^{\perp}=\{\mathbf{0}\}$.
\item $\mathbf{x}\neq\mathbf{0}$ is \hlt{isotropic} or \hlt{null} if $\innerprod{\mathbf{x},\mathbf{x}}=0$, and \hlt{nonisotropic} otherwise.\\
The space $\mathcal{V}$ is \hlt{isotropic} if such an $\mathbf{x}$ exists, and \hlt{anisotropic} otherwise.
\item $\mathbf{x}$ is \hlt{degenerate} if $\mathbf{x}\perp\mathcal{V}$. The \hlt{radical} is $\operatorname{rad}(\mathcal{V})\defeq\mathcal{V}^{\perp}$
\item $\mathcal{V}$ is \hlt{nonsingular} or \hlt{nondegenerate} if $\operatorname{rad}(\mathcal{V})=\{\mathbf{0}\}$, \hlt{singular} otherwise.
\item $\mathcal{V}$ is \hlt{totally singular} or \hlt{totally degenerate} if $\operatorname{rad}(\mathcal{V})=\mathcal{V}$.
\end{enumerate}
\end{definition}

\begin{proposition} \label{prop:metric-radical-equals-bilinear-radicals} \hlt{Metric Radical as a Bilinear Radical}\\
Let $B(\mathbf{x},\mathbf{y})=\innerprod{\mathbf{x},\mathbf{y}}$ be the form of a metric vector space $\mathcal{V}$. Then $\operatorname{rad}_{L}(B)=\operatorname{rad}_{R}(B)=\operatorname{rad}(\mathcal{V})$.\\
Thus $\mathcal{V}$ is nonsingular if and only if $B$ is perfect.
\begin{proof}
\begin{enumerate}[label=(\roman*)]
\setlength{\itemsep}{0pt}
\item By definition, $\operatorname{rad}(\mathcal{V}) =\{\mathbf{x}\in\mathcal{V}:\innerprod{\mathbf{x},\mathbf{y}}=0\ \forall\,\mathbf{y}\in\mathcal{V}\} =\operatorname{rad}_{L}(B)$
\item If $\mathbf{x}\in\operatorname{rad}_{L}(B)$, then $\innerprod{\mathbf{x},\mathbf{y}}=0$ for all $\mathbf{y}$. Since $B$ is symmetric or skew-symmetric, $\innerprod{\mathbf{y},\mathbf{x}}=\pm\innerprod{\mathbf{x},\mathbf{y}}=0$, so $\mathbf{x}\in\operatorname{rad}_{R}(B)$. The reverse inclusion is identical.
\item $B$ is perfect if and only if $\operatorname{rad}_{L}(B)=\operatorname{rad}_{R}(B)=\{\mathbf{0}\}$, if and only if $\operatorname{rad}(\mathcal{V})=\{\mathbf{0}\}$.
\end{enumerate}
\hlt{Structural Commentary}: Metric forms are the square bilinear forms whose two null directions agree.\\
\hlt{Key Insight}: Metric geometry uses the common left-right radical.
\end{proof}
\end{proposition}

\begin{proposition} \label{prop:metric-radical-representing-matrix} \hlt{Radical and Representing Matrices}\\
Let $\beta=(\mathbf{v}_{1},\ldots,\mathbf{v}_{n})$ be an ordered basis of $\mathcal{V}$.\\
Then $\{[\mathbf{x}]_{\beta}:\mathbf{x}\in\operatorname{rad}(\mathcal{V})\}=\ker(M_{\beta}^{t})$, and $\dim\operatorname{rad}(\mathcal{V})=\dim(\mathcal{V})-\operatorname{rank}(M_{\beta})$.
\begin{proof}
For $\mathbf{c}=[\mathbf{x}]_{\beta}$, $\mathbf{x}\in\operatorname{rad}(\mathcal{V})$ $\Longleftrightarrow$ $\innerprod{\mathbf{x},\mathbf{v}_{j}}=0\ \forall j$ $\Longleftrightarrow$ $\mathbf{c}^{t}M_{\beta}=0$ $\Longleftrightarrow$ $\mathbf{c}\in\ker(M_{\beta}^{t})$.\\
Rank-nullity gives $\dim\operatorname{rad}(\mathcal{V})=n-\operatorname{rank}(M_{\beta}^{t})=n-\operatorname{rank}(M_{\beta})$.\\
\hlt{Structural Commentary}: Degenerate vectors are precisely the coordinate null vectors of the representing matrix.\\
\hlt{Key Insight}: The radical is the coordinate nullspace of $M_{\beta}^{t}$.
\end{proof}
\end{proposition}

\begin{proposition} \label{prop:metric-nonsingular-matrix-invertible} \hlt{Matrix Criterion for Nonsingularity}\\
Let $\beta$ be an ordered basis of $\mathcal{V}$. Then $\mathcal{V}$ is nonsingular if and only if $M_{\beta}$ is invertible.
\begin{proof}
\begin{enumerate}[label=\roman*.]
\setlength{\itemsep}{0pt}
\item[$(\Rightarrow)$] If $\operatorname{rad}(\mathcal{V})=\{\mathbf{0}\}$, then $\ker(M_{\beta}^{t})=\{\mathbf{0}\}$. Since $M_{\beta}$ is square, $M_{\beta}$ is invertible.
\item[$(\Leftarrow)$] If $M_{\beta}$ is invertible, then $\ker(M_{\beta}^{t})=\{\mathbf{0}\}$, so $\operatorname{rad}(\mathcal{V})=\{\mathbf{0}\}$.
\end{enumerate}
\hlt{Structural Commentary}: Zero radical is matrix invertibility in coordinates.\\
\hlt{Key Insight}: Nonsingularity of a form is ordinary matrix nonsingularity in any basis.
\end{proof}
\end{proposition}

\begin{remark} \label{rem:subspace-radical-versus-ambient-complement} \hlt{Subspace Radicals}\\
For $\mathcal{S}\leq\mathcal{V}$, compute $\operatorname{rad}(\mathcal{S})$ using $\innerprod{\cdot,\cdot}|_{\mathcal{S}\times\mathcal{S}}$, where $\operatorname{rad}(\mathcal{S}) =\{\mathbf{s}\in\mathcal{S}:\innerprod{\mathbf{s},\mathbf{t}}=0\ \forall\,\mathbf{t}\in\mathcal{S}\} =\mathcal{S}\cap\mathcal{S}^{\perp}$.\\
In general, $\operatorname{rad}(\mathcal{S})\neq\mathcal{S}^{\perp}$; the latter is ambient.
\end{remark}

\begin{proposition} \label{prop:riesz-map-isomorphism-criterion} \hlt{Riesz Map Isomorphism Criterion}\\
Let $\mathcal{V}$ be a finite-dimensional metric vector space and define $\tau:\mathcal{V}\to\mathcal{V}'$ where $\tau(\mathbf{x})(\mathbf{v})=\innerprod{\mathbf{v},\mathbf{x}}$.\\
Then $\ker(\tau)=\operatorname{rad}(\mathcal{V})$, and $\tau$ is an isomorphism $\Longleftrightarrow$ $\operatorname{rad}(\mathcal{V})=\{\mathbf{0}\}$.
\begin{proof} Bilinearity gives $\tau\in\LL(\mathcal{V},\mathcal{V}')$, and\\
 $\mathbf{x}\in\ker(\tau) \Longleftrightarrow \innerprod{\mathbf{v},\mathbf{x}}=0\ \forall\,\mathbf{v}\in\mathcal{V} \Longleftrightarrow \innerprod{\mathbf{x},\mathbf{v}}=0\ \forall\,\mathbf{v}\in\mathcal{V} \Longleftrightarrow \mathbf{x}\in\mathcal{V}^{\perp} =\operatorname{rad}(\mathcal{V})$
\begin{enumerate}[label=\roman*.]
\setlength{\itemsep}{0pt}
\item[$(\Rightarrow)$] If $\tau$ is an isomorphism, then $\ker(\tau)=\{\mathbf{0}\}$, hence $\operatorname{rad}(\mathcal{V})=\{\mathbf{0}\}$.
\item[$(\Leftarrow)$] If $\operatorname{rad}(\mathcal{V})=\{\mathbf{0}\}$, then $\ker(\tau)=\{\mathbf{0}\}$. Since $\dim\mathcal{V}=\dim\mathcal{V}'$, $\tau$ is an isomorphism.
\end{enumerate}
\hlt{Structural Commentary}: The Riesz map is invertible exactly off the radical.\\
\hlt{Key Insight}: Nonsingularity is injectivity of the Riesz map.
\end{proof}
\end{proposition}

\begin{theorem} \label{thm:riesz-map-for-nonsingular-metric-spaces} \hlt{Riesz Map for Nonsingular Metric Spaces}\\
Let $\mathcal{V}$ be a finite-dimensional nonsingular metric vector space. Then $\forall f\in\mathcal{V}'\ \exists!\,\mathbf{x}\in\mathcal{V}$, $f(\mathbf{v})=\innerprod{\mathbf{v},\mathbf{x}}\ \forall\,\mathbf{v}\in\mathcal{V}$.
\begin{proof}
Let $\tau:\mathcal{V}\to\mathcal{V}'$ be given by $\tau(\mathbf{x})(\mathbf{v})=\innerprod{\mathbf{v},\mathbf{x}}$. By Proposition \ref{prop:riesz-map-isomorphism-criterion}, $\tau$ is an isomorphism.\\
Hence each $f\in\mathcal{V}'$ has a unique $\mathbf{x}\in\mathcal{V}$ with $\tau(\mathbf{x})=f$.\\
\hlt{Structural Commentary}: Nonsingularity identifies vectors with covectors; positivity is irrelevant.\\
\hlt{Key Insight}: Riesz representation is fundamentally nonsingularity, not order or length.
\end{proof}
\end{theorem}

\begin{definition} \label{def:metric-adjoint-nonsingular-form} \hlt{Adjoint for a Nonsingular Bilinear Form}\\
Let $\mathcal{V}$ be nonsingular.\\
For $T\in\LL(\mathcal{V})$, the \hlt{metric adjoint} is the unique $T^*\in\LL(\mathcal{V})$ such that $\innerprod{T\mathbf{x},\mathbf{y}}=\innerprod{\mathbf{x},T^*\mathbf{y}}$ $\forall\,\mathbf{x},\mathbf{y}\in\mathcal{V}$.
\end{definition}

\begin{proposition} \label{prop:metric-adjoints-and-isometries} \hlt{Metric Adjoints and Isometries}\\
Let $\mathcal{V}$ be nonsingular and let $T\in\LL(\mathcal{V})$.
\begin{enumerate}[label=(\roman*)]
\setlength{\itemsep}{0pt}
\item $\exists!\,T^*\in\LL(\mathcal{V})$.
\item If $\beta$ is ordered, $M=M_{\beta}$, and $A=[T]_{\beta}$, then $[T^*]_{\beta}=M^{-1}A^{t}M$.
\item $T$ is an isometry $\Longleftrightarrow$ $T$ is bijective and $T^*T=I$. For bijective $T$, this is equivalent to $T^{-1}=T^*$.
\end{enumerate}
\begin{proof}
\begin{enumerate}[label=(\roman*)]
\setlength{\itemsep}{0pt}
\item For $\mathbf{y}\in\mathcal{V}$, set $f_{\mathbf{y}}(\mathbf{x})=\innerprod{T\mathbf{x},\mathbf{y}}$. Theorem \ref{thm:riesz-map-for-nonsingular-metric-spaces} gives a unique $T^*\mathbf{y}$ such that $f_{\mathbf{y}}(\mathbf{x})=\innerprod{\mathbf{x},T^*\mathbf{y}}$ for all $\mathbf{x}$.\\
Uniqueness gives linearity of $\mathbf{y}\mapsto T^*\mathbf{y}$.
\item With $A=[T]_{\beta}$ and $B=[T^*]_{\beta}$, the adjoint identity gives $A^{t}M=MB$, so $B=M^{-1}A^{t}M$.
\item $\forall \mathbf{x},\mathbf{y}\in\mathcal{V}$, $\innerprod{T\mathbf{x},T\mathbf{y}} =\innerprod{\mathbf{x},T^*T\mathbf{y}}$. so form preservation $\Longleftrightarrow \innerprod{\mathbf{x},(T^*T-I)\mathbf{y}}=0\ \forall\,\mathbf{x},\mathbf{y}\Longleftrightarrow T^*T=I$.\\
This gives the stated bijective criterion; if $T$ is bijective, then $T^*T=I\Longleftrightarrow T^*=T^{-1}$.
\end{enumerate}
\hlt{Structural Commentary}: Adjoint theory needs nonsingularity, not positivity.\\
\hlt{Key Insight}: Isometry means the adjoint is the inverse.
\end{proof}
\end{proposition}

